% ============================================================================
% DRAFT v2 -- Section 7: Conclusion
% Genre: benchmark / evaluation study.
% ============================================================================

\section{Conclusions}
\label{sec:conclusion}

This paper belongs to the class of studies that measure how a field evaluates
itself. Such studies do not add a detector; they establish what the field's
existing measurements are worth, and they are useful to the extent that the
answer changes how the next comparison is read. Our answer is that the standard
evaluation design for unsupervised anomaly detection cannot resolve the
differences the field reports, and that in the pipeline we instrumented the
design axis those comparisons are usually about contributes almost nothing.

The binding constraint is resolution. A two-way decomposition of the 210-detector
matrix we trained leaves a residual standard deviation of $3.30$ points, and the
run-to-run standard deviation of a single detector at a fixed (backbone,
category) is $1.96$ points over three seeds. Comparing two backbones across the
fifteen categories the benchmark provides is therefore a paired test whose minimum
detectable difference is $2.98$ points --- $75\%$ of the entire $3.98$-point
spread of the fourteen backbones, and six times the $0.5$-point margin that
separates PatchCore from the strongest results reported since. Resolving that
margin on a single category would take $240$ independent training runs per
backbone, or $532$ categories. This is a statement about the benchmark rather than
about our pipeline, and it bounds every architecture comparison made on it.

Within that constraint, the diagnosis of the architecture axis is unambiguous.
Two thirds of the search space is spent on a feature dimension that cannot reach
the deployed model, because the encoder that consumes it hard-codes a width while
the decoder honours the requested one. The search signal does not distinguish
architectures: it separates the eleven tight well-behaved backbones by $5.7\%$ of
its median, which is the same order as its own run-to-run noise, and it returns
exactly $0.0$ when a student diverges --- for two backbones in most trials, and
occasionally for any of them. And the quantity it is used to predict is dominated
something else entirely: architecture accounts for $0.54\%$ of the variance in
detection accuracy and the category for $95.32\%$
(Section~\ref{subsec:capacity_null}).

The interaction between the first two facts produces the result we did not
anticipate. Because the objective can silently return the optimum of its own
range, and the search minimises a single sample per candidate, the rule does not
select the best architecture but the one with the heaviest left tail --- which is
the one that failed. Twelve independent runs selected a configuration prone to
this collapse every single time, and the architectures they returned perform
\emph{worse} than uniform random selection, at the $1.2$th percentile of the
random distribution (Section~\ref{subsec:search_vs_random}). The four runs whose
results the manuscript originally reported chose a backbone that does not
collapse, which on this account was luck.

The constructive finding is where the accuracy actually is. The reduction that
turns a multi-scale discrepancy map into an image-level score is worth $+7.92$
points of mean AUROC across three training seeds --- an order of magnitude more
than the architecture --- and it is worth those points without retraining
anything. The default choice, an unweighted mean over all pixels, is not merely
suboptimal: it fails below chance on \texttt{screw}, and top-$k$ pooling repairs
that failure to $84\%$. The gain is robust across two orders of magnitude in $k$
(Section~\ref{sec:leverage}). And the reduction can be chosen without any anomaly
label, by perturbing held-out normal images with synthetic defects and selecting
the policy that best separates them from the clean ones: averaged over three
training seeds this recovers $61\%$ of the oracle gap, though the seeds range from
$41\%$ to $72\%$, so it is a direction rather than a settled number.

Three limitations bound these claims. The diagnosis is established on MVTec-AD,
with one teacher, one detector family and one training protocol; whether the same
collapse appears in other pipelines is a question our evidence cannot settle,
though the mechanism --- a thresholded objective that can return its own optimum,
combined with a minimising selection rule --- is not specific to this codebase.
Our recommendation for the architecture axis is deliberately narrow: the backbone
effect is statistically detectable ($p = 0.039$) and merely two orders of
magnitude smaller than the category effect, so we do not conclude that
architecture does not matter, nor that NAS for anomaly detection is worthless in
principle. And the resolvability figure is conditional on this benchmark's fifteen
categories and this task's residual variance; it should be recomputed rather than
assumed for another benchmark or detector family --- which is exactly the
calculation a benchmark paper ought to publish, and which here decides the
question.

Two things follow for practice. A search whose objective can silently return the
optimum of its own range must filter degenerate evaluations before minimising, and
should report the distribution of its per-configuration scores rather than a
single minimum. And any deployed reverse-distillation detector currently scoring
images by the mean over its discrepancy map can be re-scored, on the models it
already has, with a $k$ of two percent, and expect to gain several points of
accuracy and to stop failing below chance on the categories that currently defeat
it. Neither of those requires convincing anyone that a new architecture is better,
which is convenient, because on this benchmark that cannot be shown.
